% Conversión TeX fiel del cuerpo científico del delta narrativo.
% Fuente congelada: ../../../../../HMT2/CONTINUIDAD_EDITORIAL/RECONSTRUCCION_ARMONICA_MACROTRATADO_20260822/REVISION_CONJUNTA_PARTES_I_II/auditar_pdf_rigor_academico/docs/DELTA_NARRATIVO_UNIDAD_RADION_AUTOCONSERVACION_AUTOINERCIA_GRAVEDAD_PENTADICA_20260828.md:12-366.
\subsection{Typed Separation among the Curvature Regime \(\tau\in\{+1,0,-1\}\), the Radial Response \(p\), and the Global Spin Holonomy}

The radion's biography can be represented by a raised curve.

\[
\gamma(\theta)
=
\bigl(r(\theta)\cos\theta,\,
r(\theta)\sin\theta,\,
m(\theta)\bigr),
\]

where the first two coordinates belong to the Euclidean plane \(\mathbb R^2\), and \(m(\theta)\) records the memory or depth of the lift. This equation permits a joint reading of circle, helix, and spiral.

When the angle evolves while the radius and memory remain constant, the planar publication is a circle. When angle and memory evolve while the radius retains its value, the complete trajectory is a helix of constant radius. When angle and radius evolve over a fixed memory, a spiral appears in \(\mathbb R^2\). When all three coordinates evolve, the trajectory is a radial helix: a biography in which phase, amplitude, and memory are jointly transformed.

The radial character \(p\) classifies the response law of the radius. One useful family is

\[
r_p(\theta)
=
r_0\bigl[1+p\lambda(\theta-\theta_0)\bigr]^{1/p},
\qquad p\neq 0,
\]

with logarithmic limit

\[
r_0\exp\bigl(\lambda(\theta-\theta_0)\bigr)
\qquad\text{when }p\to0.
\]

In this parameterization, \(p=2\) publishes the Fermat regime, \(p=1\) the Archimedean regime, \(p=0\) the logarithmic regime, \(p=-1\) the hyperbolic regime, and \(p=-2\) the lituus regime, with the corresponding choices of domain and orientation. Fractional values describe intermediate regimes. Thus, \(p=\tfrac12\), \(p=\tfrac32\), or \(p=-\tfrac12\) express radial responses that interpolate among the integer families and extend the available geometric classification.

The character can also change when crossing a threshold. The differential writing

\[
\frac{dx}{d\theta}
=
\beta(m)\,x^{1-p(m)},
\qquad x=\frac r{r_0},
\]

describes an evolution through regimes: \(p(m)\) remains stable within a phase and changes when the memory reaches a publication boundary. The resulting trajectory connects circular, helical, and spiral sectors through a common law. Continuity of the state and the threshold ledger make those sectors phases of a single evolution.

Three invariants occupy here distinct and coordinated levels. The curvature sign

\[
\tau\in\{+1,0,-1\}
\]

classifies the elliptic, parabolic, and hyperbolic regimes. The character \(p\) classifies the radial response. Spin records the orientation holonomy accumulated by the lift. Curvature, radial response, and spin form a reading system: curvature determines the local type of space, \(p\) determines how the trajectory occupies that space, and spin preserves the orientation acquired while traversing it. Closed and open curves then appear as distinct geometric biographies of the same oriented unit.

\subsection{Particle as a Persistent Section of an Enriched Orbit: Phase, Route, Fibre, Multisection, and Memory}

Dimensional Mechanics can formulate the particle as a persistent section of that geometric and dynamical extension. A compact conceptual notation is

\[
\mathfrak P
=
\operatorname{Sec}_{\mathrm{pers}}
\bigl[
\tau,p,\operatorname{Hol};
\text{sheet},\text{fiber},\text{route},\text{memory}
\bigr].
\]

The particle preserves a recognizable route through changes of phase, curvature, and dimension. Its spin publishes orientation holonomy; its charge publishes phase polarity; its color publishes an internal orientation of the gauge fiber; its flavor distinguishes families of routes or multisections; its mass expresses the dimensional-transduction character by which areal action acquires a volumetric reading. Each property is a reader of the persistent state, and all refer to a common genealogy.

The electron occupies the basal threshold of this construction because it offers the first stable reader in which phase, charge, spin and mass are simultaneously published. The rector electronic mass

\[
m_e^\beta=0.510998950690000808608\ldots\ \mathrm{MeV}
\]

fixes a transduction reference; the preceding basal states preserve their function as genealogical stages. From the electron onward, the quotients \(K_D/K_\Omega\), the towers \(90/120\), and the domain \(81/56/13/324\) organize the publication of massive families. The mass law therefore reads how each persistent section traverses the enriched space and converts areal density of action into volumetric presence.

This interpretation connects spirals to particles precisely. The index \(p\) supplies the radial response; \(\tau\) supplies the curvature regime; holonomy supplies spin; the route and the fiber supply flavor and color; area--volume transduction supplies the mass character. The particle atlas thereby becomes a classification of stable modes of dimensional traversal.

\subsection{Transduction of the oriented phase in electric polarization, magnetic response, and boundary radiative emission}

Electricity is presented as directed transport of phase and action through a vacancy structure. A vacancy is an admissible transition of the support: it opens a channel, fixes a threshold, and permits the update to advance. The electric component publishes transport tangential to the route. The magnetic component publishes the transverse curvature induced by that transport. Their perpendicularity expresses the geometric relation between advance and rotation.

When two sinusoidal transverse components maintain the appropriate phase difference, the field vector rotates. Longitudinal propagation lifts that rotation to a helix. Variation of amplitude transforms its frontal projection into a spiral. The geometric reading is then fully coordinated: angle carries phase; angular velocity carries frequency; radius carries amplitude and action density; orientation carries polarization and direction; the vertical coordinate carries propagation and memory; \(p\) carries the radial regime; the threshold carries emission, absorption, or change of state.

Electromagnetism appears as a phase–action–memory dynamics with complementary publications. Electricity records polar advance; magnetism records the transverse response; the wave records periodicity; the helix records advance with memory; the spiral records radial evolution; radiation records the publication of the state when accumulation reaches a threshold.

This architecture also orders the role of vacancies and cycles. Vacancy supplies the local possibility of transition; the cycle preserves phase; memory distinguishes turns; the threshold determines the next residence of the action. Propagation may remain confined, change radial regime, or be published as controlled radiation. Field forces thereby acquire a continuous genealogical description from the discrete support to the observed wave.

\subsection{Reversibility of the complete dynamics and thermodynamic cost of information-reducing projections}
The Landauer zero represents a limiting case of this law. For a bijective update of the complete state,

\[
H(X\mid U(X))=0,
\qquad
\inf Q_L=0.
\]

The information remains available in the genealogy. When a reader forms a quotient \(q\), entropy decomposes as

\[
H(X)=H(q(X))+H(X\mid q(X)).
\]

The conditional term measures the information retained outside the projection. A materialized tritic reset satisfies the threshold

\[
Q_{\mathrm{reset}}\ge k_B\Theta\log 3.
\]

The Boltzmann constant translates the multiplicity grouped by the thermal reader; the Nyquist threshold separates faithful sampling from aliasing; null Landauer characterizes the reversible transport of the integral state. Information, temperature and action are articulated by the type of reading applied.

\subsection{Self-conservative forces as transport of action among the observable,
radial, memory, vacancy, and radiation sectors}

Evolutionary conservation is better expressed through exchange than through scalar immobility. A complete state distributes action and momentum among visible motion, memory, radial deformation, vacancies, radiation, and the conjugate return branch. The enriched balance may be written as

\[
\Delta J_{\mathrm{visible}}
+\Delta J_{\mathrm{memoria}}
+\Delta J_{\mathrm{radial}}
+\Delta J_{\mathrm{vacancias}}
+\Delta J_{\text{radiation}}
+\Delta J_{\mathrm{retorno}}
=0.
\]

A force is \textbf{self-conserving} when its own update law contains the channels that transport, store, publish, and reabsorb action. Conservation belongs to the complete genealogical balance of the interaction. Each reader observes part of the exchange; the genealogy reconstructs its totality.

The Earth--Moon system provides a direct physical image. Earth's rotation, the lunar orbit, tides, deformation, and radiation exchange angular momentum. The net consumption of the complete system closes in the genealogical balance of those transfers. The magnitude called energy functions as a reading derived from the action and momentum published in each sector. Fundamental conservation resides in compatibility among all readers of the exchange.

This formulation dissolves the former boundary between conservative forces and so-called nonconservative forces through a deeper category. Friction, apparent dissipation, and radiation occupy explicit channels of the extended state. A loss in the visible mechanical projection corresponds to a transfer toward memory, deformation, heat, vacancies, or radiation. Self-conservation follows the residence of action through those changes.

\subsection{Self-inertial systems: publication of acceleration through orientation,
curvature, sheet, and memory}

The appropriate expression is \textbf{relational self-inertia}. It denotes the compatibility between the total state's own connection and the response that appears when it is projected within the global distribution. The endogenous and exogenous aspects are analytic components of a single dynamic relation.

Let \(\Gamma\) be a horizontal trajectory or geodesic of the complete state. Its intrinsic law satisfies

\[
\nabla_{\dot\Gamma}\dot\Gamma=0.
\]

When it is projected onto a reader \(S\), the observed acceleration is determined by the curvature of the projection:

\[
\nabla^S_{\dot\Gamma_S}\dot\Gamma_S
=
\mathcal K_{\pi_S}(\dot\Gamma,\dot\Gamma).
\]

The published acceleration expresses the relation between total motion and the geometry of the reader. Mach's principle completes this mechanism through the global boundary trace

\[
b:\mathcal X_\infty\longrightarrow B_\infty
\]

and the factorization of the local inertial response

\[
I_v=\bar I_v\circ b.
\]

The global distribution enters the local response through a typed map.
Relational self-inertia therefore combines the endogenous connection of
motion, the global condition transmitted by \(b\), and the
observer's concrete projection. A synthetic expression for this
coordination is

\[
\nabla^{\mathrm{auto}}
=
\mathcal M\bigl(
\nabla^{\mathrm{end}},
b[\mathcal X_\infty],
\pi_S
\bigr),
\]

where \(\mathcal M\) represents the Machian coupling and preserves the type of each component.

The Ambivalence Principle adds the double orientation of reading. Every direct publication preserves a conjugate return branch; emission and absorption form related endpoints of the same biography; temporal orientation is read from both directions of transport. Absorber theory finds a natural formulation here: source and receiver jointly close the genealogical balance of phase and action. The torsional present is the interface at which both orientations are rendered compatible.

TRIT acts at this level as a time crystal: it organizes the local periodicity of updating and, through memory, preserves the difference between successive returns. The direct branch and the conjugate branch relate emission to absorption, and the observer--observable coupling acquires a dynamical residence within the same state. Published time is the reading of that orientation; genealogical time is the complete sequence of phase, return, and memory.

Mach and Ambivalence perform complementary functions. Mach relates the local response to the global distribution. Ambivalence preserves the forward and conjugate directions of evolution. Relational self-inertia articulates both and replaces the division between inertial and non-inertial systems with a single law of connection, projection, and memory.

\subsection{Area--volume duality, hexad--octad transition, and the role of the Barbero--Immirzi parameter}

The passage from area to volume constitutes a central dimensional
transduction. Area publishes boundary action; volume publishes interior
depth. The HMT derivation of the Barbero–Immirzi parameter organizes the
compatibility between the two measures through hexad–octad incidence and
the channels \(90/120\):

\[
90=6\binom62,
\qquad
120=8\binom62.
\]

The cardinals \(6\) and \(8\) admit a cubic realization by means of faces and vertices, whereas the Steiner hexad and octad are distinct incidence objects linked only by the typed morphism \(6\to8\). The \(12\) edges of the cube retain their own combinatorial function and do not by themselves define Steiner incidence. The difference in incidence expresses a productive dimensional incompatibility: surface information requires a transduction rule in order to acquire volumetric residence.

The same principle appears in the carry \(999+1=1000\). The nines reach the boundary of the decimal register; the increment simultaneously publishes the closure of the preceding level and the origin \(0/1\) of the next. The holographic ratio among \(729\), \(243\), \(27\), and \(1\) unfolds that transition at triadic scales. The Barbero–Immirzi transduction is inserted into this family of reader changes: boundary and interior preserve the unit while redistributing its measure.

The normalized defect \(-1/12\) and the oriented weight \(12:-1\) perform distinct functions within this geometry. The former quantifies an incidence correction; the latter records a signed orientation in the transduction. Their coordination with the minimal area, the modular Hamiltonian, and memory provides the bridge to dimensional dynamics.

\subsection{Hyperbolic interior contraction, exterior expansion, and holographic conservation of the boundary}

The thesis may be formulated from a single idea: unity preserves its identity while changing scale, reader, curvature, memory, and dimension. This permanence has an evolutionary character because unity unfolds, acquires new coordinates, and publishes different properties; it has a holographic character because each finite publication retains the rule for reconstructing the state from which it proceeds. The Holographic Closure of Infinity designates precisely this compatibility between identity and unfolding: each stage contains the law that permits the same genealogy to be deepened inward and prolonged outward.

The governing designation is \textbf{principle of evolutionary conservation of the unit through dimensional transduction and holographic projection}. The expression places each term at its proper level. Conservation pertains to genealogical identity; evolution pertains to updating; transduction relates dimensional realizations; holographic projection preserves the reconstructive correspondence between its readers.

Elementary arithmetic offers the first image of this architecture. The decimal completion

\[
1000=10^3=729+243+27+1=3^6+3^5+3^3+3^0
\]

exposes four triadic strata within a decimal cubic unit. Upon normalization,

\[
1=\frac{729}{1000}+\frac{243}{1000}+\frac{27}{1000}+\frac{1}{1000},
\]

the unit is recognized simultaneously as totality and as graded unfolding. The reciprocal reading

\[
3^{-6},\quad 3^{-5},\quad 3^{-3},\quad 3^0
\]

describes the inner depth of those strata; its normalization produces the corresponding partition. Relation \(999+1=1000\) expresses the transport mechanism: the last unit completes a register and acts simultaneously as the carry inaugurating the next. Zero and one thus form a continuity interface. The intervals \([0,1]\), \([0,10]\), and the prolongation toward \([0,\infty)\) are scalar readers of a single completion law.

The first splitting \(1000=729+271\), followed by \(271=243+27+1\), shapes the extreme and mean holographic ratio of the deployment: a dominant triadic core preserves the same resolution rule within the remainder. The decimal cubic unit thus contains a nested triadic depth. Movement toward \(729\), \(243\), \(27\), and \(1\) reads the inward contraction; the carry toward one thousand reads the outward expansion.

This double orientation organizes HMT infinity. Inward, refinement grows hyperbolically: each cell admits a new depth without losing its filiation. Outward, unity expands through successive dimensional publications. Closure consists in the commutativity of both movements. If \(U\) represents the update of the state, \(R_d\) the reader on stratum \(d\), and \(T_{d\to d'}\) the transduction between strata, evolutionary conservation is expressed by means of

\[
R_{d'}\circ U
=
T_{d\to d'}\circ R_d.
\]

The same event may provide an angular measure, an action density, a published mass, or a gravitational response. Their common identity resides in the genealogy that makes those readings commute.

\subsubsection{Nonadic Monodromy and Helical Lifting with Memory Advance}

APP, TRIT, and TPK provide the formal mechanism for this unfolding. APP constructs the support and separates the evaluations that must be preserved. TRIT organizes the local regimes and their ambivalence. TPK transports phase and memory over the enriched support. The causal dependence is fixed by

\[
U_t\longrightarrow \Gamma_9\longrightarrow K_{\mathrm{ph}}.
\]

Phase return and memory advance form coordinated operations. After nine steps, the same angular class may reappear,

\[
g(K+9)=g(K),
\]

as the memory counter progresses,

\[
q_{\mathrm{mem}}\Gamma_9=q_{\mathrm{mem}}+1,
\qquad
B_{K+9}\neq B_K.
\]

The precise geometric figure is a helical lift. A closed orbit in the angular projection lifts to an open trajectory when each phase return increments the memory. The circle preserves the cycle; the helix preserves the cycle and additionally records the depth reached. The coinductive monodromy thereby generates infinite depth through a sequence of locally closed and globally ascending returns.

This image has more than merely illustrative value. The angular wheel constitutes a section of the helix; the helix constitutes the wheel's complete biography when the state transports memory. Each turn recovers a phase position and simultaneously initiates a new sheet. Infinitude arises from compatibility between return and lifting: iterating the cycle produces depth without dissolving orbital identity.

The nonadic certificate gives finite and verifiable content to that architecture. The chain

\[
w_6\longrightarrow w_{12}\longrightarrow w_{18}\longrightarrow
w_{24}\longrightarrow w_{30}\longrightarrow R_{36}\longrightarrow G_9
\]

orders the lift; the \(396\) levels, the \(2376\) trits, the \(43\) turns, and the \(387\) pairs \((K,K+9)\) per channel quantify the transit; the five regions of \(\pi\) and the ambient fiber of \(729\) elements publish subsequent internal readings. This sequence makes a productive monodromy visible: the return recovers the phase and coinduction preserves the growth of memory.

\subsubsection{Rational unit and coordination of the circular, elliptic, projective, and helical readers}

The radion concentrates within a metrological unit the structure that the radian publishes angularly. Its natural definition is that of an \textbf{oriented section of phase, action, and memory}. One unit of phase \(\theta=1\) transports one unit of action \(\sigma\hbar_{\mathrm{ret}}\), belongs to a complete turn \(T_+=2\pi_{\mathrm{HMT}}\), and preserves sheet, orientation, nonadic return, and position in the refinement cylinder.

The radian expresses the visible angular coordinate. The radion preserves, together with that coordinate, the transported action, orientation, carry, and memory of the revolution. Angular metrology thereby acquires a genealogy. Measuring an arc means identifying how much rotation has been published and which complete state has realized that rotation.

The relationship between phase and area takes on a particularly clear form:

\[
\mathcal A(\theta)=\hbar_{\mathrm{ret}}\,\theta.
\]

One radional unit sweeps an action area \(\hbar_{\mathrm{ret}}\); one complete turn encloses \(h_{\mathrm{ret}}\). The passage from phase to area is incorporated into the reading unit itself. The angular metric, action, and memory cease to appear as juxtaposed quantities and become coordinates of the same oriented section.

The radion admits several coordinated geometric publications. The circular reader publishes phase; the elliptical reader publishes action and form; the Klein interior reader publishes depth; the nonadic helical reader publishes return and memory. The unit remains genealogically identical across those realizations. The radion therefore constitutes the appropriate metrological link between HMT and Dimensional Mechanics: it offers a unit that can traverse changes of reader while preserving the state presupposed by each partial measurement.

\subsection{Five-Component Gravitational Polarity and Joint Transduction of Action, Area, Volume, Time, and Mass}

Dimensional Mechanics makes it possible to define dimension as a stable publication regime characterized by the minimum rank of independent memory required by the lift. To traverse a dimension is to apply a transducer that preserves genealogy and changes the type of reading. Gravity occupies the level of connection between these regimes.

In this formulation, \textbf{interdimensional gravitation with five correlative polarizations}, abbreviated as the \textbf{pentapolar structure of gravitation}, is the connection that transports genealogical unity and leaves minimal torsion at each interface. A compact notation for this trace is

\[
\Theta_d^{\min}
=
\nabla_{d+1}\circ T_{d\to d+1}
-
T_{d\to d+1}\circ\nabla_d.
\]

The expression measures the difference between transporting after differentiating and differentiating after transporting. That difference is the minimal geometric memory of the dimensional crossing. The composition of interfaces produces a gravitational holonomy

\[
\mathfrak G_{d\to D}
=
\operatorname{Hol}
\bigl(
\Theta_d^{\min},
\Theta_{d+1}^{\min},
\ldots,
\Theta_{D-1}^{\min}
\bigr).
\]

The fundamental object is this interdimensional connection. A local graviton-like mode occupies the derived rank of the quantized publication of torsion; gravitation retains its primary residence in compatibility among dimensions.

The term pentapolar encompasses the five polarizations inherent to the continuous state:

\[
\mathbf G_d^{\pm}
=
\bigl(
G_{\mathrm{sp},d}^{\pm},
G_{\mathrm{sol},d}^{\pm},
G_{\mathrm{gau},d}^{\pm},
G_{\mathrm{coh},d}^{\pm},
G_{\mathrm{det},d}^{\pm}
\bigr).
\]

The five components emerge correlatively from the same enriched state. The polarity \(\pm\) records direct concentration and conjugate publication. Gravity traverses the dimensional strata while simultaneously transporting those five readings; minimal torsion preserves the trace of each crossing.

Geometry thereby acquires a unified image. The elliptic regime organizes closures; the parabolic regime organizes thresholds; the hyperbolic regime organizes depth. Dimensional transduction connects these regimes with the radial responses \(p\), spin holonomies, and mass publications. The pentapolar structure of gravitation expresses the global coherence of that connection.

M-theory provides subsequent physical screens for this connection: additional dimensions, strings, branes, and T-duality realize concrete modes of transduction. HMT supplies the APP–TRIT–TPK genealogy, the enriched state, memory, and bidirectional reading; the interface with M-theory publishes those structures in typed physical geometries. T-duality exchanges compatible dimensional readers and preserves the genealogical class of the state.

\subsubsection{Compatibility of the classical, relativistic, quantum, and HMT-MD realizations}

Classical mechanics publishes trajectory and momentum. Relativity publishes causal structure, metric, and curvature. Quantum mechanics publishes phase, holonomy, vacancies, and amplitudes. Dimensional Mechanics provides the transducers that relate those publications, and HMT preserves the discrete genealogy that sustains them.

In the classical reading, an orbit describes the visible projection of motion. In the relativistic reading, that orbit is incorporated into the geometry of the reader. In the quantum reading, phase and spin record the holonomy of the section. In the HMT--MD reading, the three expressions belong to the same lifted biography: the radion quantifies phase, action, and memory; \(p\) classifies the radial response; \(\tau\) classifies curvature; holonomy publishes spin; the dimensional transducer publishes mass; the pentapolar connection publishes gravitation.

Apparent classical dissipation acquires a residence in the expanded genealogical balance. Relativistic acceleration is expressed through connection and projection. Quantum indeterminacy is organized by readers, vacancies, and branches of Ambivalence. Mach's principle relates the local response to the global distribution. Zero Landauer cost guarantees the informational reversibility of the complete state. All these elements converge in a deeper conservation: the identity of the unit across changes of publication.

\subsubsection{Genealogical Conservation and Holographic Prolongation Narrative Theorem}

The closure can now be narrated as a single sequence. The unit decomposes without losing identity. The oriented phase acquires measure in the radion. The nonadic return converts the cycle into a helix through memory. The character \(p\) transforms the helix into radial families and classifies their regimes. The persistence of those trajectories produces sections with spin, charge, flavour, colour, and mass. Vacancies open channels; electrical polarity transports phase; magnetic response curves transversely; thresholds publish radiation. Self-conservation closes the genealogical balance of action and momentum. Relational self-inertia links the proper connection, projection, and Machian totality. Ambivalence preserves the direct and conjugate directions. Area--volume transduction leads to Barbero--Immirzi and the mass law. Minimal torsion links dimensions and publishes the pentapolar structure of gravitation.

Inward growth and outward expansion are conjugate orientations of the same process. Hyperbolic refinement generates depth; dimensional transduction generates extension. Carry \(0/1\) links scales; monodromy links turns; memory links states; holonomy links orientations; gravitation links dimensions. Each linkage preserves an evolving unity, and every finite publication contains the rule of its prolongation.

The fundamental thesis is formulated as follows:

\begin{quote}
\textbf{The holographic closure of infinity realizes the principle of the evolutionary conservation of unity through dimensional transduction and holographic projection: unity maintains its genealogical identity as it unfolds as phase, action, curvature, memory, particle, field, and gravitation; every finite reading preserves the rule of reconstruction and prolongation of the complete state.}
\end{quote}

This formulation provides a joint physical and mathematical image. Infinity acquires a law of local closure; the unit acquires depth; dimension acquires memory; conservation acquires relational character; force acquires self-preservation; inertia acquires Machian structure; gravity acquires a pentapolar connection; and holography acquires the precise meaning of reconstruction among compatible readers.
